% ===================================================================
%  TABLEAU N° 8 — La récolte.
%
%  Traduction, faite en 2026, en francais standard du Quebec, sur
%  l'ido de texto/io. Regles de la colonne : voir 10-tableau-01.tex.
%  Deux scenes, deux numerotations : les cles portent c1 et c2.
%
%  LES SIX CHOIX PROPRES A CE TABLEAU :
%    (1)  bluets ............... centaurées
%    (10) pâtre ................ berger
%    (13) berger ............... gardien du troupeau
%    (47) machine à battre ..... batteuse
%    (10) carnier .............. gibecière
%    (53) bicyclette ........... vélo
%
%  « BLUET » EST LE PIEGE LE PLUS SERIEUX DE TOUTE LA COLONNE, parce
%  qu'il ne se voit pas. Le fac-simile ecrit « Bluets, coquelicots et
%  digitales » : ce sont les fleurs des champs, et « bluet » y nomme
%  la Centaurea cyanus, bleue, qui pousse dans les bles. Au Quebec, le
%  bleuet est un FRUIT — la baie du Lac-Saint-Jean — et rien d'autre.
%  Le lecteur d'ici aurait lu des petits fruits melanges a des
%  coquelicots au milieu d'un champ de ble. « Centaurée » est le nom
%  que le Quebec donne a la fleur, et il ne prete a rien.
%
%  (10) ET (13) SONT DEUX HOMMES, ET LE FAC-SIMILE LES NOMME « pâtre »
%  ET « berger ». Le premier mot ne se dit plus ; mais on ne peut pas
%  se contenter de le remplacer par le second, qui est deja pris.
%  L'ido distingue « pastoro », le berger, et « muton-gardisto », qui
%  garde les moutons sous ses ordres — celui qui boit a la gourde
%  pendant que l'autre surveille l'orage. D'ou berger et gardien du
%  troupeau. Le berger du tableau 7 (renvoi 45) est le meme mot ido et
%  porte deja le meme nom.
% ===================================================================

%%K t08-noto-1 noto
\VUnotes{143.2mm}{%
\textit{(*) Parce que ce mot n'est pas précis : s'agit-il de la récolte du lin,
de celle du raisin, de celle des pommes? Il serait peut-être bon d'employer
«} \VUgras{mesono} \textit{», proposé dans Mondo XIV, p. 242. Mais jusqu'ici
cette racine nouvelle n'a pas été adoptée par l'Académie et attend d'être
discutée selon la procédure habituelle des idistes.}%
}

%%K t08-apar-1 apar
\VUsaut{5.0mm}
\VUtitre{15.0pt}{60}{TABLEAU N\textsuperscript{o} 8}
\VUsaut{3.0mm}
\VUcentre{13.2pt}{90}{\textit{Première scène.}}
\VUsaut{3.0mm}
\VUpk{10.2pt}{40}{La récolte (*).}
\VUsaut{4.0mm}
\VUpk{10.2pt}{40}{Les aspects de la campagne.}
\VUsaut{3.0mm}

%%K t08-c1-01-1 p
1. --- Avec l'\VUgras{été}, l'aspect de la campagne a changé une fois de plus.
La semence confiée au sillon a germé; une petite plante verte est sortie du sol
et a monté en épi. Le grain s'est formé, puis il a mûri, et dès maintenant il
n'y a plus qu'à récolter.

%%K t08-c1-02-1 p
2. --- Les \VUgras{fleurs des champs} sont épanouies. \VUgras{Centaurées}
\textsuperscript{(1)}, \VUgras{coquelicots} \textsuperscript{(2)} et
\VUgras{digitales} \textsuperscript{(3)} mêlent à la teinte uniforme des champs
de blé le gai contraste de leurs \VUgras{corolles} fraîches.

%%K t08-c1-03-1 p
3. --- Les arbres fruitiers se sont garnis de feuilles; le fruit a mûri. Le
petit \VUgras{cerisier} \textsuperscript{(4)} est couvert de belles
\VUgras{cerises} \textsuperscript{(5)} que les enfants cueillent et mangent avec
grand plaisir.

%%K t08-c1-04-1 p
4. --- Désormais, le bétail peut passer toute la journée en plein air.

%%K t08-c1-04-2 p
Les \VUgras{agneaux} \textsuperscript{(6)} ont grandi et sont devenus de belles
\VUgras{brebis} \textsuperscript{(7)} et de beaux \VUgras{moutons}
\textsuperscript{(8)} à l'épaisse toison. Ils broutent paisiblement
l'\VUgras{herbe} du \VUgras{pâturage} \textsuperscript{(9)} émaillé de
\VUgras{marguerites}.

%%K t08-c1-04-3 p
Bientôt on va tondre \VUgras{ces} bêtes pour avoir leur \VUgras{laine}, dont on
fait le drap de nos vêtements.

%%K t08-tit-2 sub
\VUsaut{5.0mm}
\VUpk{10.2pt}{40}{Le berger.}
\VUsaut{3.4mm}

%%K t08-c1-05-1 p
5. --- Le \VUgras{berger} \textsuperscript{(10)} semble ne pas beaucoup faire
attention à eux. Il ne se soucie que de l'orage qui passe à l'horizon, et le
\VUgras{gardien du troupeau} \textsuperscript{(13)}, qui approche sa
\VUgras{gourde} \textsuperscript{(14)} de ses lèvres, a l'air encore plus
insouciant.

%%K t08-c1-06-1 p
6. --- La chaleur est accablante, parce que le \VUgras{chien}
\textsuperscript{(15)} vient lui aussi se désaltérer; il lape l'eau fraîche du
\VUgras{ruisseau} \textsuperscript{(16)} et effraie une pauvre petite
\VUgras{grenouille} \textsuperscript{(17)} qui s'était tapie dans les herbes de
la rive. Elle saute dans l'eau claire où fleurissent les \VUgras{nénuphars}
\textsuperscript{(18)}.

%%K t08-c1-07-1 p
7. --- Une \VUgras{bergeronnette} \textsuperscript{(19)} se baigne près de la
petite \VUgras{source} \textsuperscript{(20)} qui jaillit entre deux rochers et
se cache sous les \VUgras{broussailles} \textsuperscript{(21)} et les
\VUgras{buissons d'aubépine} \textsuperscript{(22)}.

%%K t08-tit-3 sub
\VUsaut{5.0mm}
\VUpk{10.2pt}{40}{La récolte.}
\VUsaut{3.4mm}

%%K t08-c1-08-1 p
8. --- Pour les enfants de la campagne, la récolte du blé est une époque bien
amusante. Ils font de joyeuses \VUgras{culbutes} \textsuperscript{(24)} sur les
\VUgras{tas de paille} \textsuperscript{(23)}, ils aident les
\VUgras{moissonneurs} \textsuperscript{(26)} et, pendant que ceux-ci fauchent le
\VUgras{champ de blé} \textsuperscript{(27)} avec leurs \VUgras{faux}
\textsuperscript{(25)} bien aiguisées sur la \VUgras{pierre à aiguiser}
\textsuperscript{(30)}, ils ramassent le blé et le lient en \VUgras{gerbes}
\textsuperscript{(31)} ou en \VUgras{javelles} \textsuperscript{(32)}.

%%K t08-c1-09-1 p
9. --- Parfois, on permet aux plus grands d'entre eux de couper avec une
\VUgras{faucille} \textsuperscript{(12)} les \VUgras{tiges dorées}
\textsuperscript{(28)} qui portent les lourds \VUgras{épis}
\textsuperscript{(29)} pleins de grain; et les petits, tout en gardant le
\VUgras{bétail} \textsuperscript{(33)} qui broute tranquillement dans le
\VUgras{pâturage} \textsuperscript{(34)}, à l'ombre du grand \VUgras{tilleul}
\textsuperscript{(35)}, attrapent dans leur \VUgras{filet}
\textsuperscript{(36)} les beaux \VUgras{papillons}.

%%K t08-c1-09-2 p
Quelques-uns grimpent même sur la \VUgras{charrette} \textsuperscript{(37)}
pour ranger les gerbes qu'on leur passe au bout d'une \VUgras{fourche}
\textsuperscript{(38)}.

%%K t08-c1-10-1 p
10. --- Dans toute la \VUgras{plaine} \textsuperscript{(39)}, où s'envolent les
\VUgras{alouettes} \textsuperscript{(41)}, règne l'animation. Tout le monde est
occupé à la récolte. Mais, pendant que les pauvres continuent d'employer les
vieux instruments, \VUgras{faux}, \VUgras{faucilles} et \VUgras{fléaux}
\textsuperscript{(40)}, les riches propriétaires font faucher leur blé ou leur
\VUgras{seigle} \textsuperscript{(43)} avec une \VUgras{moissonneuse mécanique}
\textsuperscript{(42)}. Ensuite, comme il n'est pas nécessaire de porter les
gerbes à la grange, on forme sur place de grandes \VUgras{meules}
\textsuperscript{(44)}.

%%K t08-c1-11-1 p
11. --- On amène alors une \VUgras{batteuse} \textsuperscript{(47)} qu'une
\VUgras{locomobile} \textsuperscript{(45)} met en mouvement à l'aide d'une
\VUgras{courroie de transmission} \textsuperscript{(46)}, et le grain se sépare
ainsi de la \VUgras{paille} sans quitter le champ où il a été semé.

%%K t08-tit-4 sub
\VUsaut{5.0mm}
\VUpk{10.2pt}{40}{L'orage.}
\VUsaut{3.4mm}

%%K t08-c1-12-1 p
12. --- Il arrive souvent de violents \VUgras{orages} \textsuperscript{(53)} au
temps de la récolte. On entend d'abord au loin les grondements du
\VUgras{tonnerre}; un violent \VUgras{vent d'ouest} \textsuperscript{(54)} se
met à souffler et courbe en haletant les plus gros arbres du \VUgras{bois}
\textsuperscript{(49)}. De gros \VUgras{nuages} \textsuperscript{(56)} noirs
envahissent le ciel; ils sont traversés par les zigzags aveuglants de
l'\VUgras{éclair} \textsuperscript{(55)}. La \VUgras{pluie} et parfois la
\VUgras{grêle} se mettent à tomber, écrasant les récoltes prêtes à être
fauchées.

%%K t08-c1-13-1 p
13. --- Souvent la foudre met le feu à une construction. C'est ainsi que l'an
dernier le toit du \VUgras{moulin à vent} \textsuperscript{(50)} a été réduit en
cendres et que deux de ses \VUgras{ailes} \textsuperscript{(51)} ont été
fracassées.

%%K t08-c1-14-1 p
14. --- Après quelques heures, le calme renaît. Un brillant \VUgras{arc-en-ciel}
\textsuperscript{(52)} apparaît à l'horizon et la nature reprend sa vie,
interrompue quelques instants. Les gens de la campagne, qui s'étaient réfugiés
dans les \VUgras{cabanes} \textsuperscript{(48)}, se remettent au travail et
s'efforcent courageusement de réparer les dégâts qu'ils viennent de subir.

%%K t08-tit-5 sub
\VUsaut{6.0mm}
\VUcentre{13.2pt}{90}{\textit{Deuxième scène.}}
\VUsaut{3.6mm}

%%K t08-tit-6 sub
\VUpk{10.2pt}{40}{La chasse. --- Les vendanges.}
\VUsaut{1.4mm}
\VUpk{10.2pt}{40}{La pêche.}
\VUsaut{4.0mm}

%%K t08-c2-01-1 p
1. --- Vers la fin d'\VUgras{août} ou le milieu de \VUgras{septembre}, selon les
régions, c'est l'ouverture de la chasse. À peine les premiers rayons du soleil
ont-ils fait sortir les \VUgras{lézards} \textsuperscript{(2)} de leur trou que
le \VUgras{chasseur} \textsuperscript{(5)} se met en route avec ses
\VUgras{chiens} \textsuperscript{(4)}.

%%K t08-c2-02-1 p
2. --- Il a mis son solide \VUgras{costume de chasse} \textsuperscript{(9)} de
gros drap, il a chaussé des \VUgras{guêtres} de cuir qui lui permettront de
marcher dans l'herbe mouillée où se cachent les \VUgras{crapauds}
\textsuperscript{(1)}; il a pris son \VUgras{fusil} \textsuperscript{(6)} et sa
\VUgras{gibecière} \textsuperscript{(10)}, et le voilà parti.

%%K t08-c2-03-1 p
3. --- L'ardeur de la poursuite l'amène au milieu d'un vignoble où les vendanges
vont bon train. Les enfants du vigneron, qui gardent d'ordinaire les chèvres au
bord du chemin, les laissent aujourd'hui brouter à l'aventure et, après avoir
attaché à un pieu un vieux \VUgras{bouc} \textsuperscript{(3)} qui ne manquerait
pas de profiter de sa liberté pour s'échapper, s'accordent à leur tour une part
du plaisir. L'un saisit une grappe de raisin dans un \VUgras{panier à vendange}
\textsuperscript{(12)} et s'amuse à faire couler le \VUgras{jus}
\textsuperscript{(14)} dans un \VUgras{gobelet} \textsuperscript{(13)} pour
faire du moût (du vin qui n'a pas fermenté). L'autre se coiffe d'une
\VUgras{couronne de feuillage} \textsuperscript{(15)} qu'il vient de tresser. À
côté d'eux, des \VUgras{moineaux} \textsuperscript{(16)} sans gêne viennent
jusque devant le pressoir picorer les raisins.

%%K t08-c2-04-1 p
4. --- Pendant que les enfants s'amusent, les hommes travaillent. Les
\VUgras{vendangeurs} \textsuperscript{(18)} portent le raisin au cellier dans
des \VUgras{hottes} \textsuperscript{(17)}; un \VUgras{tonnelier}
\textsuperscript{(19)} répare les \VUgras{barriques} \textsuperscript{(20)},
vérifie si les \VUgras{douves} \textsuperscript{(22)} et les \VUgras{fonds}
\textsuperscript{(21)} sont en bon état et remplace les \VUgras{cercles}
\textsuperscript{(23)} cassés. C'est lui aussi qui a fait les grandes
\VUgras{cuves} \textsuperscript{(25)} dans lesquelles on verse le
\VUgras{raisin} \textsuperscript{(26)} pour le faire fermenter.

%%K t08-c2-05-1 p
5. --- Quand la fermentation est terminée, on tire le vin et on met le résidu
sur le \VUgras{pressoir} \textsuperscript{(28)} et, en serrant progressivement
la \VUgras{vis} \textsuperscript{(29)}, on exprime le jus, qui coule dans un
\VUgras{baquet} \textsuperscript{(32)}, et l'on obtient encore du \VUgras{vin}
\textsuperscript{(31)}.

%%K t08-tit-8 sub
\VUsaut{6.0mm}
\VUpk{10.2pt}{40}{Bière et cidre.}
\VUsaut{3.4mm}

%%K t08-c2-06-1 p
6. --- Sur le mur du \VUgras{cellier} \textsuperscript{(27)} grimpe une branche
de \VUgras{houblon} \textsuperscript{(30)}. Ce n'est là qu'une plante
décorative, mais dans certains pays, où la vigne est très rare, on cultive le
houblon pour faire de la \VUgras{bière}.

%%K t08-c2-07-1 p
7. --- Le petit \VUgras{pommier} \textsuperscript{(33)} est chargé de pommes qui
donneront, quand elles seront mûres, un excellent \VUgras{cidre}. Dans leur
\VUgras{cage} \textsuperscript{(34)}, le \VUgras{serin} \textsuperscript{(35)}
et le \VUgras{chardonneret} \textsuperscript{(36)} regardent d'un œil envieux
les moineaux qui folâtrent en liberté.

%%K t08-c2-08-1 p
8. --- À perte de vue, sur le versant des collines, s'alignent les
\VUgras{échalas} \textsuperscript{(40)} qui portent les magnifiques \VUgras{ceps
de vigne} \textsuperscript{(39)}, chargés de lourdes \VUgras{grappes}.

%%K t08-c2-08-2 p
Pendant toute l'année, le \VUgras{vigneron} \textsuperscript{(42)} a travaillé
pour soigner son \VUgras{vignoble} \textsuperscript{(38)}, et le voilà
récompensé de sa peine par une belle récolte. Les \VUgras{vendangeuses}
\textsuperscript{(41)} remplissent les cuves posées sur la charrette tirée par
des \VUgras{mules} \textsuperscript{(43)}, et tout le monde est joyeux.

%%K t08-tit-9 sub
\VUsaut{5.0mm}
\VUpk{10.2pt}{40}{La pêche.}
\VUsaut{3.4mm}

%%K t08-c2-09-1 p
9. --- Il en est de même des \VUgras{pêcheurs à la ligne} \textsuperscript{(44)}
qui, depuis le matin, sont venus s'installer au bord de l'eau avec leurs
\VUgras{cannes à pêche} \textsuperscript{(45)}.

%%K t08-c2-09-2 p
Le \VUgras{poisson} \textsuperscript{(47)} mord à l'\VUgras{hameçon} dès qu'ils
jettent leur \VUgras{ligne} \textsuperscript{(46)} à l'eau, et ils ont à peine
le temps de renouveler l'\VUgras{appât} qu'une nouvelle victime frétille au bout
de la ligne.

%%K t08-c2-09-3 p
Pourtant, le \VUgras{pêcheur au filet} \textsuperscript{(48)} qui, monté dans sa
\VUgras{barque} \textsuperscript{(49)}, lance l'\VUgras{épervier} au milieu de
la \VUgras{rivière} \textsuperscript{(50)} prend plus de poissons.

%%K t08-c2-10-1 p
10. --- L'automne est la saison des belles promenades : à pied, à
\VUgras{cheval} \textsuperscript{(52)}, à \VUgras{vélo} \textsuperscript{(53)},
en \VUgras{automobile} \textsuperscript{(54)}. Les \VUgras{routes}
\textsuperscript{(51)} sont propres et l'on profite des derniers beaux jours,
parce que voilà déjà les \VUgras{hirondelles} \textsuperscript{(56)} qui se
rassemblent sur les toits pour s'envoler à tire-d'aile vers des pays plus
chauds. Le feuillage du vieux \VUgras{noyer} \textsuperscript{(57)} jaunit, les
\VUgras{noix} tombent, et les grands \VUgras{arbres} groupés en
\VUgras{bouquet} \textsuperscript{(58)} sur la \VUgras{hauteur}
\textsuperscript{(59)} seront bientôt sans feuilles.

%%K t08-c2-11-1 p
11. --- Le \VUgras{soleil} \textsuperscript{(60)} se lève plus tard, les jours
raccourcissent, un petit froid assez piquant commence à se faire sentir le
\VUgras{matin}, et voilà déjà des vols de \VUgras{bécasses}
\textsuperscript{(61)} qui quittent nos climats devenus inhospitaliers.
\VUsaut{6.0mm}
\VUfilet{22mm}
